\begin{explanationbox}
	\textbf{\textcolor{CExplanation}{一句话直觉}}

	求导公式就是一把钥匙，将每个初等函数与它的变化率直接挂钩.

	\medskip
	\textbf{\textcolor{CExplanation}{核心拆解}}
	\begin{description}[style=nextline, leftmargin=2.8em, labelsep=0.5em, itemsep=0.45em]
		\item[\textbf{要点一（幂函数导数）：}] 在适用范围内，幂函数求导是“指数前移、指数减一”。根式与倒数可先改写为幂，但不能扩大定义域。
		\item[\textbf{要点二（指对函数导数）：}] 指数函数保留原函数乘以$\ln a$，对数函数导数是$1/(x\ln a)$. $e^x$和$\ln x$是最简洁的特例.
		\item[\textbf{要点三（三角导数）：}] $\sin$ 变 $\cos$，$\cos$ 变 $-\sin$；正负号不可混淆.
	\end{description}

	\medskip
	\textbf{\textcolor{CExplanation}{几何本质（切线斜率）}}
% F049: native TikZ/PGFPlots, computed from the stated mathematical model.
\begin{center}
\begin{tikzpicture}[font=\small,>=stealth]
\begin{axis}[width=.90\linewidth,height=5.5cm,xmin=-0.3,xmax=4.8,ymin=-0.3,ymax=2.7,
axis lines=left,xlabel={$x$},ylabel={$y$},font=\small,
xlabel style={at={(axis description cs:1.02,.02)},anchor=west},
ylabel style={rotate=-90,at={(axis description cs:.02,1.03)},anchor=south},
legend style={at={(axis description cs:.5,-.18)},anchor=north,draw=none,fill=none,font=\small,legend columns=2},
unbounded coords=jump,clip=true,]
\addplot[gray!45,thin,densely dotted,forget plot] coordinates {(-0.3,0) (4.8,0)};
\addplot[gray!45,thin,densely dotted,forget plot] coordinates {(0,-0.3) (0,2.7)};
\addplot[blue!70!black,thick] coordinates {(0,0) (0.025,0.15811388) (0.05,0.2236068) (0.075,0.27386128) (0.1,0.31622777) (0.125,0.35355339) (0.15,0.38729833) (0.175,0.41833001) (0.2,0.4472136) (0.225,0.47434165) (0.25,0.5) (0.275,0.52440442) (0.3,0.54772256) (0.325,0.57008771) (0.35,0.59160798) (0.375,0.61237244) (0.4,0.63245553) (0.425,0.65192024) (0.45,0.67082039) (0.475,0.68920244) (0.5,0.70710678) (0.525,0.72456884) (0.55,0.74161985) (0.575,0.75828754) (0.6,0.77459667) (0.625,0.79056942) (0.65,0.80622577) (0.675,0.82158384) (0.7,0.83666003) (0.725,0.85146932) (0.75,0.8660254) (0.775,0.88034084) (0.8,0.89442719) (0.825,0.90829511) (0.85,0.92195445) (0.875,0.93541435) (0.9,0.9486833) (0.925,0.9617692) (0.95,0.97467943) (0.975,0.98742088) (1,1) (1.025,1.0124228) (1.05,1.0246951) (1.075,1.0368221) (1.1,1.0488088) (1.125,1.0606602) (1.15,1.0723805) (1.175,1.0839742) (1.2,1.0954451) (1.225,1.1067972) (1.25,1.118034) (1.275,1.129159) (1.3,1.1401754) (1.325,1.1510864) (1.35,1.161895) (1.375,1.1726039) (1.4,1.183216) (1.425,1.1937336) (1.45,1.2041595) (1.475,1.2144958) (1.5,1.2247449) (1.525,1.2349089) (1.55,1.24499) (1.575,1.25499) (1.6,1.2649111) (1.625,1.2747549) (1.65,1.2845233) (1.675,1.2942179) (1.7,1.3038405) (1.725,1.3133926) (1.75,1.3228757) (1.775,1.3322913) (1.8,1.3416408) (1.825,1.3509256) (1.85,1.3601471) (1.875,1.3693064) (1.9,1.3784049) (1.925,1.3874437) (1.95,1.396424) (1.975,1.4053469) (2,1.4142136) (2.025,1.4230249) (2.05,1.4317821) (2.075,1.440486) (2.1,1.4491377) (2.125,1.457738) (2.15,1.4662878) (2.175,1.4747881) (2.2,1.4832397) (2.225,1.4916434) (2.25,1.5) (2.275,1.5083103) (2.3,1.5165751) (2.325,1.5247951) (2.35,1.532971) (2.375,1.5411035) (2.4,1.5491933) (2.425,1.5572412) (2.45,1.5652476) (2.475,1.5732133) (2.5,1.5811388) (2.525,1.5890249) (2.55,1.5968719) (2.575,1.6046807) (2.6,1.6124515) (2.625,1.6201852) (2.65,1.6278821) (2.675,1.6355427) (2.7,1.6431677) (2.725,1.6507574) (2.75,1.6583124) (2.775,1.6658331) (2.8,1.6733201) (2.825,1.6807736) (2.85,1.6881943) (2.875,1.6955825) (2.9,1.7029386) (2.925,1.7102631) (2.95,1.7175564) (2.975,1.7248188) (3,1.7320508) (3.025,1.7392527) (3.05,1.7464249) (3.075,1.7535678) (3.1,1.7606817) (3.125,1.767767) (3.15,1.7748239) (3.175,1.781853) (3.2,1.7888544) (3.225,1.7958285) (3.25,1.8027756) (3.275,1.8096961) (3.3,1.8165902) (3.325,1.8234583) (3.35,1.8303005) (3.375,1.8371173) (3.4,1.8439089) (3.425,1.8506756) (3.45,1.8574176) (3.475,1.8641352) (3.5,1.8708287) (3.525,1.8774983) (3.55,1.8841444) (3.575,1.890767) (3.6,1.8973666) (3.625,1.9039433) (3.65,1.9104973) (3.675,1.917029) (3.7,1.9235384) (3.725,1.9300259) (3.75,1.9364917) (3.775,1.9429359) (3.8,1.9493589) (3.825,1.9557607) (3.85,1.9621417) (3.875,1.968502) (3.9,1.9748418) (3.925,1.9811613) (3.95,1.9874607) (3.975,1.9937402) (4,2) (4.025,2.0062403) (4.05,2.0124612) (4.075,2.0186629) (4.1,2.0248457) (4.125,2.0310096) (4.15,2.0371549) (4.175,2.0432817) (4.2,2.0493902) (4.225,2.0554805) (4.25,2.0615528) (4.275,2.0676073) (4.3,2.0736441) (4.325,2.0796634) (4.35,2.0856654) (4.375,2.0916501) (4.4,2.0976177) (4.425,2.1035684) (4.45,2.1095023) (4.475,2.1154196) (4.5,2.1213203)};
\addlegendentry{$f(x)=\sqrt x$}
\addplot[orange!90!black,thick] coordinates {(0.286,0.643) (1.714,1.357)};
\addlegendentry{$x_0=1,\ k=0.5$}
\addplot[only marks,mark=*,orange!90!black,forget plot] coordinates {(1,1)};
\addplot[teal!75!black,thick,dashed] coordinates {(0.847,0.9325886) (1.765,1.337057)};
\addlegendentry{割线：$\Delta y/h$}
\addplot[gray,densely dotted,forget plot] coordinates {(1,1) (1.612,1) (1.612,1.2696456)};

\end{axis}
\end{tikzpicture}
\par\smallskip{\small 绿色割线的斜率是差商；另一点趋近切点时，割线趋向橙色切线。}
\end{center}



	导数$f'(x_0)$就是曲线$y=f(x)$在点$(x_0,f(x_0))$处切线的斜率. 因此公式直接给出各函数图像的切线斜率.

	\medskip
	\textbf{\textcolor{CExplanation}{代数意义（变化率）}}

	导数刻画函数值随自变量变化的瞬时快慢. 例如$(\ln x)'=1/x$说明$x$越大增长越慢.

	\medskip
	\textbf{\textcolor{CExplanation}{考点价值（求导运算基础）}}
	\begin{itemize}[leftmargin=*, itemsep=0.5em, topsep=0.4em]
		\item \textbf{考法一（直接求导）：} 给出函数解析式，要求计算导函数. 必须准确套用公式.
		\item \textbf{考法二（切线问题）：} 求曲线在某点或过某点的切线方程，必须先用导数公式算出斜率.
	\end{itemize}

	\medskip
	\textbf{\textcolor{CExplanation}{顿悟点}}

	\textbf{公式中的每个细节（底数、指数、符号）都对应一个运算规则，记熟了求导就是“照方抓药”.}

	\medskip
	\textbf{\textcolor{CExplanation}{使用场景}}
	\begin{itemize}[leftmargin=*, itemsep=0.5em, topsep=0.4em]
		\item \textbf{场景一（基础求导）：} 由基本初等函数及其线性组合构成的函数，直接逐项求导.
		\item \textbf{场景二（切线斜率计算）：} 给定切点横坐标，代入导函数得斜率，再写切线方程.
	\end{itemize}

\medskip
\textbf{\textcolor{CExplanation}{新增公式如何从基本公式长出来？}}
\begin{description}[style=nextline,leftmargin=2.8em,labelsep=0.5em,itemsep=0.45em]
\item[\textbf{根式和倒数：}] $\sqrt{x}=x^{1/2}$（$x>0$），$1/x=x^{-1}$（$x\ne0$）。先认出幂函数，再求导。
\item[\textbf{正切与余切：}] $\tan x=\sin x/\cos x$，$\cot x=\cos x/\sin x$。商法则的分子由 $\sin^2x+\cos^2x=1$ 化简；正切导数为正，余切导数为负。
\item[\textbf{绝对值：}] $y=|x|$ 在原点左右分别是斜率 $-1$ 和 $1$ 的直线，两个斜率不相等，所以原点不可导。
\item[\textbf{变量既在底数又在指数：}] $x^x=e^{x\ln x}$，先把变量指数移到指数函数的内层，再用乘积及复合求导。
\item[\textbf{反三角函数（拓展）：}] 反函数交换横纵坐标，相应切线斜率互为倒数（原斜率须非零）。
\end{description}
\medskip
\textbf{\textcolor{CExplanation}{速记关联：}} $\tan$ 的导数为 $1+\tan^2x$；$\cot$ 的导数为 $-(1+\cot^2x)$。$\ln|x|$ 左右两支导数均是 $1/x$。

\medskip
\textbf{\textcolor{CExplanation}{核心证明学习顺序：}} 先理解两个基本极限（指对基本极限与三角基本极限），再学习七条公式的差商变形。指数提公因式、对数化商、实数幂转指数、三角和差化积，是四类不同的化简思路。
\end{explanationbox}
